\documentclass[10pt]{article}
\usepackage[margin=1in]{geometry}
\usepackage[T1]{fontenc}
\usepackage{lmodern,microtype}
\usepackage{amsmath,amssymb,amsthm}
\usepackage[hidelinks]{hyperref}
\newtheorem{theorem}{Theorem}
\newtheorem{lemma}[theorem]{Lemma}
\newcommand{\F}{\mathbf F}
\newcommand{\Q}{\mathbf Q}
\newcommand{\R}{\mathbf R}
\newcommand{\lk}{\operatorname{lk}}
\newcommand{\GL}{\operatorname{GL}}
\newcommand{\PGL}{\operatorname{PGL}}
\newcommand{\PG}{\operatorname{PG}}
\title{A negative answer to Nevo's\\homology-sphere decomposition question}
\author{Alec Kriebel\\
\small Independent researcher\\
\small \href{https://orcid.org/0009-0001-9320-500X}{ORCID: 0009-0001-9320-500X}}
\date{17 October 2026 \quad (version 1.0.0)}
\begin{document}
\maketitle
\begin{abstract}
We construct a finite two-dimensional abstract simplicial complex that is doubly
Cohen--Macaulay over every field of characteristic zero, but contains no
two-dimensional homology-sphere subcomplex over any field. This gives a negative
answer to Nevo's question asking whether every doubly Cohen--Macaulay complex
admits a decomposition into homology spheres with specified overlap properties.
We use a finite quotient of an affine building, with all links equal to the
incidence graph of the projective plane of order eleven. A direct spectral
estimate controls the links after any vertex deletion; a weighted cocycle
argument verifies the required homology vanishing. The sphere obstruction is
an Euler-characteristic count. The construction and the spectral and curvature
mechanisms build on established work, credited below.
\end{abstract}

\section{The question and the result}
All complexes in this note are finite abstract simplicial complexes unless a
building is explicitly being discussed. A pure complex $K$ is Cohen--Macaulay
over a field $k$ if, for every face $\sigma$, including the empty face,
\[
 \widetilde H_i(\lk_K\sigma;k)=0
 \qquad\text{for all } i<\dim(\lk_K\sigma).
\]
It is doubly Cohen--Macaulay, or $2$-CM, if it is Cohen--Macaulay and the deletion
$K-v$ of every vertex is Cohen--Macaulay of the same dimension.

Nevo proved that a $d$-dimensional $2$-CM complex can be expressed as a union of
minimal $d$-cycle complexes with overlaps containing a $d$-face and a
reordering condition, for $d\geq1$ \cite[Theorem~7]{NevoOWR}
\cite[Theorem~3.4]{Nevo}. He then asked whether the pieces could instead be
homology spheres \cite[Problem~8, p.~961]{NevoOWR}
\cite[Problem~3.7, p.~417]{Nevo}. We use the standard local-link convention for
a homology sphere: a $d$-dimensional homology sphere $S$ over $k$ satisfies
\[
 \widetilde H_i(\lk_S\sigma;k)=
 \begin{cases}
 k,&i=d-|\sigma|,\\
 0,&i\ne d-|\sigma|
 \end{cases}
 \quad\text{for every face }\sigma,
\]
including $\sigma=\varnothing$. Here $\{\varnothing\}$ is the sphere of
dimension $-1$. This convention, also used in the surrounding sphere and ear
decomposition literature, is explicitly stated in \cite[Section~4]{Swartz}.
Having the homology of a sphere only globally is a weaker condition.

\begin{theorem}\label{thm:main}
There exists a finite pure two-dimensional abstract simplicial complex $K$
which is $2$-CM over every characteristic-zero field and has no
two-dimensional homology-sphere subcomplex over any field.
Consequently Nevo's homology-sphere decomposition question has a negative answer.
\end{theorem}

A single such complex over $\Q$ disproves the universal question. The finite
field $\F_{11}$ used in the geometry is distinct from the homology coefficient
field; no claim of $2$-CM over every positive-characteristic field is made.
No subdivision, immersion, or enlarged ambient complex is used.

The substantial ingredients are classical. Lubotzky--Samuels--Vishne
\cite[Sections~2--4]{LSV} provide the affine building, a vertex-regular lattice
and a faithful polynomial matrix representation. The method comparing polynomial
degree with congruence divisibility to obtain large displacement already appears in
Lubotzky--Meshulam \cite[Proposition~3.3]{LM}; we give the particular matrix
reduction needed here. The link spectral vanishing criterion is a
Garland-type result, for example \cite[Theorem~1]{BS}; we include its elementary
finite two-dimensional proof. The curvature obstruction to triangulated
two-spheres with links of length at least six is recorded in
\cite[Example~1.8(5)]{JS}. The point of this note is their application to the
precise $2$-CM sphere question, with every vertex-deletion condition verified.

\section{A finite quotient with projective-plane links}
Let $F=\F_{11}((y))$ and $\mathcal O=\F_{11}[[y]]$. The affine building
$\mathcal B$ of $\PGL_3(F)$ is two-dimensional. Its vertices are homothety
classes of $\mathcal O$-lattices in $F^3$; chains of intermediate lattices define
its faces. At a vertex represented by $L$, the neighbors correspond to
nonzero proper subspaces of $L/yL\cong\F_{11}^3$. An edge in its link is an
incident point--line pair. Thus every vertex link is the bipartite incidence
graph $G$ of $\PG(2,11)$.

We recall the exact features of \cite[Sections~3--4]{LSV} used here. Their
Cartwright--Steger group $\Gamma\leq\PGL_3(F)$ acts simply transitively on the
vertices of $\mathcal B$. One description starts with the nine-dimensional
cyclic algebra over $\F_{11}(y)$ generated by $E=\F_{11^3}$ and $z$, with
\[
 za=a^{11}z\quad(a\in E),\qquad z^3=1+y.
\]
This algebra splits over $F$. The group is generated by the projective classes
of $b_u=u(1-z^{-1})u^{-1}$, for $u\in E^\times/\F_{11}^{\times}$.
Conjugation on the algebra gives a faithful representation
\[
 \rho:\Gamma\longrightarrow\GL_9(\F_{11}[t]),\qquad t=1/y.
\]
The kernel before taking projective classes consists of central elements,
so this representation is faithful on $\Gamma$. In the basis used in
\cite[Equation~(9) and Definition~4.7]{LSV}, the matrices $\rho(b_u)$ have degree
at most one. Corollary~4.6 and Proposition~4.8 there show that each neighbor
of a fixed base vertex $o$ is $b_u o$ or $b_u b_{u'}o$ for some generators.
This covers both lattice-index types, $11$ and $11^2$, and hence every
undirected neighbor. All group matrices and their inverses lie in the stated
general linear group over the polynomial ring. We reduce this faithful matrix
representation, rather than attempting to define the original cyclic algebra
over $\F_{11}[t]$. None of the later global Ramanujan spectral assertions of
\cite{LSV} is used.

Define the normal finite-index subgroup
\[
 \Gamma_9=\ker\!\left(\Gamma\xrightarrow{\rho}
 \GL_9(\F_{11}[t])\longrightarrow
 \GL_9(\F_{11}[t]/(t^9))\right).
\]
If $d(o,\gamma o)\leq4$ in the building's edge graph, follow a path and use simple
transitivity at its consecutive vertices. Each step is represented by a product
of at most two generators, so $\rho(\gamma)$ has degree at most eight.
For $\gamma\in\Gamma_9$, every entry of $\rho(\gamma)-I$ is divisible by $t^9$.
The degree bound forces $\rho(\gamma)=I$, and faithfulness forces $\gamma=1$.
Normality and vertex transitivity therefore imply
\[
 d(x,\gamma x)>4
 \quad\text{for every vertex }x
 \text{ and every }1\ne\gamma\in\Gamma_9.                 \tag{1}
\]
This is a direct polynomial argument for the nonreduced quotient ring
$\F_{11}[t]/(t^9)$. It does not invoke the irreducible-polynomial hypotheses
in the different arithmetic quotient of \cite[Proposition~3.3]{LM}.

Let $K$ have as vertices the $\Gamma_9$-orbits of building vertices, and as faces
exactly the images of building faces. We add no further faces by completing
cliques. The vertex set is finite because $\Gamma_9$ has finite index in the
vertex-regular group. By (1), distinct vertices in any face remain distinct,
and distinct neighbors of a vertex $x$ have distinct images: their distance
is at most two. Every quotient edge or triangle incident to $[x]$ lifts to
one incident to $x$, by translating a chosen lift with an element of $\Gamma_9$.
Its other vertices are uniquely determined among the neighbors of $x$.
It follows both on vertices and on edges that
\[
                   \lk_K([x])\cong G.                 \tag{2}
\]
The quotient is a connected pure two-dimensional abstract simplicial complex.

\section{Cohen--Macaulayness after every vertex deletion}
For a finite graph $H$ with positive degrees, write
\[
 \mathcal E_H(f)=\sum_{\{u,v\}\in E(H)}(f(u)-f(v))^2,
 \qquad
 \lambda(H)=\inf_{f\ {\rm nonconstant}}
 \frac{\mathcal E_H(f)}{\min_c\sum_v d_H(v)(f(v)-c)^2}.
\]
For a connected graph this is the first positive normalized Laplacian
eigenvalue.

The graph $G$ has $133$ points and $133$ lines, and is $12$-regular. If $M$ is its
point--line incidence matrix, uniqueness of the line through two distinct
points gives
\[
                      MM^{\mathsf T}=11I+J.
\]
Thus its adjacency eigenvalues are $\pm12$ and $\pm\sqrt{11}$; it is connected,
and its ordinary Laplacian gap is $12-\sqrt{11}$. Bipartiteness and uniqueness
of a line through two points exclude cycles of length less than six.

For a vertex $w$ of $G$, put $H=G-w$. Its degrees are eleven or twelve.
Extend any unweighted-mean-zero function $x$ on $H$ by zero at $w$.
The extension remains mean zero on $G$, and removing the incident edges costs
at most $\|x\|^2$. Consequently
\[
 \mathcal E_H(x)\geq(11-\sqrt{11})\|x\|^2.              \tag{3}
\]
This also proves that $H$ is connected. For an arbitrary nonconstant $f$, let
$\bar f$ be its unweighted mean on $H$. Then
\[
 \min_c\sum_v d_H(v)(f(v)-c)^2
 \leq \sum_v d_H(v)(f(v)-\bar f)^2
 \leq 12\|f-\bar f\|^2 .
\]
Together with (3) this proves
\[
 \lambda(G-w)\geq\frac{11-\sqrt{11}}{12}>\frac{7}{12}>
 \frac12,\qquad
 \lambda(G)=\frac{12-\sqrt{11}}{12}>\frac12.             \tag{4}
\]
The distinction between ordinary mean and degree-weighted variance is
necessary for the irregular graph $G-w$.

\begin{lemma}\label{lem:vanish}
Let $X$ be a finite pure two-dimensional simplicial complex. If all vertex
links are connected and have normalized Laplacian gap greater than $1/2$,
then $H^1(X;\R)=0$.
\end{lemma}
\begin{proof}
For an edge $e$, let $m(e)>0$ be the number of incident triangles.
Represent a real cohomology class by an antisymmetric edge cocycle $a$ of
minimum squared norm
$E(a)=\sum_{\{u,w\}}m(uw)a_{uw}^2$ among its representatives.
Such a minimum exists by orthogonal projection onto the finite-dimensional
space of coboundaries. Variation by a coboundary supported at $v$ gives
\[
                        \sum_u m(vu)a_{vu}=0.
\]
The function $f_v(u)=a_{vu}$ on $\lk_X v$ therefore has weighted mean zero,
since $d_{\lk_X v}(u)=m(vu)$.
For a triangle $vuw$, the cocycle equation gives
$f_v(u)-f_v(w)=-a_{uw}$. If $\lambda>1/2$ is the minimum of the finitely many
link gaps, the link Poincare inequalities, summed over $v$, give
\[
 \begin{split}
 E(a)
 &=\sum_v\sum_{\{u,w\}\in E(\lk_X v)}a_{uw}^2\\
 &\geq\lambda\sum_v\sum_{u\in V(\lk_X v)}m(vu)a_{vu}^2
 =2\lambda E(a).
 \end{split}
\]
Thus $E(a)=0$, and the original cohomology class is zero.
\end{proof}

We apply the lemma to $K$ and to every deletion $K-z$.
Each edge of $K$ belongs to twelve triangles by (2). A surviving edge loses at
most one triangle when $z$ is removed, so it still belongs to at least eleven.
Each surviving vertex has at least $265$ neighbors, hence lies in a triangle.
Thus $K-z$ remains nonempty, pure and two-dimensional.
It is connected: a path in $K$ that passes through $z$ can bypass $z$ along
its connected link $G$, whose edges are surviving edges of $K$.
For $v\ne z$, the link of $v$ in $K-z$ is $G$ if $vz$ was not an edge, and is
$G-w$ for the corresponding neighbor $w$ otherwise.
Equations (2)--(4) and Lemma~\ref{lem:vanish} give $H^1(X;\R)=0$ for
$X=K$ or $K-z$. Finite-dimensional duality gives $H_1(X;\R)=0$.
The boundary matrices have integer entries, whose ranks agree over $\Q$,
$\R$, and every characteristic-zero field. Hence the same vanishing holds
over every such field.

For completeness, these facts check all faces in the Cohen--Macaulay
criterion. At the empty face, connectedness and $H_1=0$ give all the required
below-top reduced homology vanishing. Vertex links are nonempty connected
one-dimensional complexes. Edge links are nonempty zero-dimensional complexes.
Facet links are $\{\varnothing\}$ of dimension $-1$, with no required
below-top nonzero homology. Nonemptiness covers the negative-degree
conventions. This proves Cohen--Macaulayness of $K$ and all $K-z$, of the same
dimension, and therefore the $2$-CM assertion of Theorem~\ref{thm:main}.

\section{There are no homology-sphere pieces}
Suppose that a two-dimensional homology sphere $S$ over some field is a
subcomplex of $K$. Its edge links have the reduced homology of $S^0$, so
every edge belongs to exactly two triangles of $S$. Every vertex link in $S$
is connected, and each of its vertices has degree two by that edge condition.
It is therefore a simple cycle. It embeds as a subgraph of the corresponding
link $G$ in $K$, so its length is at least six. Writing $f_i=f_i(S)$, we obtain
\[
                 2f_1\geq6f_0,\qquad 3f_2=2f_1,
 \qquad
                 \chi(S)=f_0-\frac{f_1}{3}\leq0.
\]
Global sphere homology, over any field, instead gives $\chi(S)=2$.
This contradiction proves the sphere exclusion. Neither orientability nor
a classification of surfaces is needed, and $S$ need not be induced.
In particular $K$ cannot even be covered by permitted sphere subcomplexes;
the further overlap and reordering requirements cannot be met. This completes
the proof of Theorem~\ref{thm:main}.

\section*{Verification, priority, and review disclosure}
The accompanying source and verification archive contains the paper source,
authored audit records, and reproducible finite checks. The default
standard-library verifier constructs the projective-plane incidence graph,
checks its exact Gram identity and all single-vertex deletions, and checks
a triangulated torus at the strict $1/2$ spectral boundary using rational
homology and exact cocycle energies. Actual corrupted constructions are
required to fail in normal and optimized Python modes.
A separate optional NumPy diagnostic checks the explicit polynomial
neighbor matrices and local adjacency; it is not required by the offline
default verifier. None of these programs enumerates the finite quotient,
replaces the symbolic proof, or independently constructs building theory.
Historical review programs and later packaged replays are distinguished
in the archive.

An extensive bounded priority search, including three separate adversarial
literature reviews, did not verify an earlier recorded full answer to this
specific question. This is not a certificate of firstness or a claim that
all literature has been inspected. The construction and the spectral,
congruence, and curvature tools have substantial antecedents, as credited
above. Recent work on related algebraic and numerical consequences, including
\cite{LarsonStapledon}, does not itself assert the sphere-subcomplex
decomposition considered here. Material access gaps remain for the full
bodies of Chari's 1997 convex-ear paper and Baclawski's 1982 paper; some other
original bodies, books, theses, versions, and subscription indexes were
uninspected. The exact primary versions, selective reading scopes, query
inputs, near matches, and gaps are recorded in the authored audit files.
Third-party source bodies and extracts are not redistributed.

AI tools were used extensively in developing and checking the argument,
writing, literature research, computation, and adversarial review.
This preprint is unrefereed. Internal reviews by AI agents do not constitute
independent external human peer review, journal refereeing, or formal proof
verification.

\begin{thebibliography}{9}
\bibitem{NevoOWR}
E. Nevo, contribution to \emph{Discrete Geometry}, Oberwolfach Reports
17/2005, pp.~959--962; Problem~8 on p.~961.
\href{https://ems.press/content/serial-article-files/45990}{Original workshop report}.

\bibitem{Nevo}
E. Nevo, Rigidity and the lower bound theorem for doubly Cohen--Macaulay
complexes, \emph{Discrete Comput. Geom.} \textbf{39} (2008), 411--418.
\href{https://doi.org/10.1007/s00454-007-9017-y}{doi:10.1007/s00454-007-9017-y}.
ArXiv \href{https://arxiv.org/abs/math/0505334v2}{math/0505334v2}
and the published question were inspected; Problem~3.7 is on published p.~417.

\bibitem{Swartz}
E. Swartz, $g$-elements, finite buildings and higher Cohen--Macaulay
connectivity, \emph{J. Combin. Theory Ser. A} \textbf{113} (2006), 1305--1320.
Inspected version: \href{https://arxiv.org/abs/math/0512086v1}{math/0512086v1}.

\bibitem{LSV}
A. Lubotzky, B. Samuels and U. Vishne, Explicit constructions of Ramanujan
complexes of type $\widetilde A_d$, \emph{European J. Combin.} \textbf{26}
(2005), 965--993.
\href{https://doi.org/10.1016/j.ejc.2004.06.007}{doi:10.1016/j.ejc.2004.06.007}.
Inspected version: \href{https://arxiv.org/abs/math/0406217v2}{math/0406217v2}.

\bibitem{LM}
A. Lubotzky and R. Meshulam, A Moore bound for simplicial complexes,
\emph{Bull. London Math. Soc.} \textbf{39} (2007), 353--358.
\href{https://doi.org/10.1112/blms/bdm003}{doi:10.1112/blms/bdm003}.
Proposition~3.3 was checked in the ten-page author version.

\bibitem{BS}
W. Ballmann and J. \'Swi\k{a}tkowski, On $L^2$-cohomology and property~(T)
for automorphism groups of polyhedral cell complexes,
\emph{Geom. Funct. Anal.} \textbf{7} (1997), 615--645.
Inspected version: MPI preprint 96-127 (1996),
\href{https://archive.mpim-bonn.mpg.de/id/eprint/2255/}{primary preprint record}.

\bibitem{JS}
T. Januszkiewicz and J. \'Swi\k{a}tkowski, Simplicial nonpositive curvature,
\emph{Publ. Math. Inst. Hautes \'Etudes Sci.} \textbf{104} (2006), 1--85.
\href{https://www.numdam.org/article/PMIHES_2006__104__1_0.pdf}{Published article}.

\bibitem{LarsonStapledon}
M. Larson and A. Stapledon, Complementary vectors of simplicial complexes,
\emph{Forum Math. Sigma} \textbf{14} (2026), e133.
\href{https://doi.org/10.1017/fms.2026.10288}{doi:10.1017/fms.2026.10288}.
Inspected versions: \href{https://arxiv.org/abs/2504.20264v2}{arXiv:2504.20264v2}
and the published introduction, with additional scopes recorded in the archive.
\end{thebibliography}
\end{document}
